% common/preamble.tex — paquetes compartidos por los documentos del curso.
%
% Deliberadamente ligero: el avance es un documento corto de quiz, no una
% tesis. Se carga desde la carpeta de cada documento (avance/, final/), así que
% TODAS las rutas de abajo son relativas al directorio del documento:
%   common/   -> ../common/
%   assets/   -> ../assets/pics/
%   bib/      -> ../bib/refs.bib
% ACOPLAMIENTO: si algo de eso se mueve, la compilación falla (o peor: compila
% sin el logo). No hay script que lo detecte.

% --- Matemáticas (antes de unicode-math, que si no avisa) ---
\usepackage{mathtools}

% --- Fuentes (XeLaTeX) ---
% Libertinus. Se cargan por NOMBRE DE ARCHIVO + Extension, no por nombre de
% familia: en TeXLive las .otf viven en el árbol de texmf y fontconfig no las
% ve, así que \setmainfont{Libertinus Serif} falla.
\usepackage[warnings-off={mathtools-colon,mathtools-overbracket},mathbf=sym]{unicode-math}
\setmainfont{LibertinusSerif}[
  Extension     = .otf,
  UprightFont   = *-Regular,
  ItalicFont    = *-Italic,
  BoldFont      = *-Semibold,
  BoldItalicFont= *-SemiboldItalic,
  Ligatures     = TeX,
]
\setsansfont{LibertinusSans}[
  Extension     = .otf,
  UprightFont   = *-Regular,
  ItalicFont    = *-Italic,
  BoldFont      = *-Bold,
  BoldItalicFont= *-Bold,
  Ligatures     = TeX,
]
\setmathfont{LibertinusMath-Regular.otf}

% --- Idioma ---
% es-noshorthands: los shorthands < > ~ . de babel-spanish rompen cleveref y las
% columnas d{..} de dcolumn.
\usepackage[english,spanish,es-noshorthands]{babel}
\selectlanguage{spanish}

% --- Geometría y espaciado ---
\usepackage[letterpaper,margin=2.5cm]{geometry}
\usepackage{setspace}\onehalfspacing
\usepackage{parskip}

% --- Tablas (los documentos son tabla-pesados: consignas y métricas) ---
\usepackage{booktabs}
\usepackage{longtable}
\usepackage{array}
\usepackage{tabularx}
\usepackage[section]{placeins} % las tablas no pueden flotar fuera de su sección
\newcolumntype{L}[1]{>{\raggedright\arraybackslash}p{#1}}
% Filas menos pegadas: un poco más de aire vertical en todas las tablas.
\renewcommand{\arraystretch}{1.3}
% \sepfilas: separador punteado gris discreto entre filas (booktabs no trae
% líneas punteadas). Va solo tras el \\ de cada fila de cuerpo. colortbl va
% antes de arydshln y le da el \arrayrulecolor (global: se restaura a negro).
\usepackage{colortbl}
\usepackage{arydshln} % último de los paquetes de tabla: su manual lo exige
\newcommand{\sepfilas}{%
  \arrayrulecolor{gray!60}\hdashline[0.8pt/3pt]\arrayrulecolor{black}}

% --- Color, cajas y enlaces ---
\usepackage{xcolor}
\usepackage[most]{tcolorbox}
\usepackage{enumitem}
\usepackage[hidelinks]{hyperref}

% --- Logo institucional: compartido por todos los documentos del curso ---
\usepackage{graphicx}
\graphicspath{{../assets/pics/}}

% --- Comillas dependientes del idioma (biblatex lo exige junto con babel) ---
\usepackage{csquotes}

% --- Bibliografía: UN SOLO refs.bib para todos los entregables del curso ---
\usepackage[backend=biber,style=authoryear,sorting=nyt,maxcitenames=2,
            maxbibnames=99,giveninits=true,uniquename=init,doi=true,url=true,isbn=false]{biblatex}
\addbibresource{../bib/refs.bib}

% common/macros.tex — dos macros de ayuda, compartidos por todos los documentos
% del curso (avance, final).

% ---------------------------------------------------------------------------
% Modo FINAL: `make es FINAL=1` crea el archivo .final-mode y las cajas de ayuda
% desaparecen. El flag lo busca ESTE archivo con \IfFileExists, así que es
% relativo al directorio desde donde compila latexmk (la carpeta del documento).
% ---------------------------------------------------------------------------
\newif\ifdocFinal \docFinalfalse
\IfFileExists{.final-mode}{\docFinaltrue}{}

% ---------------------------------------------------------------------------
% \limite{N}{qué pide la consigna aquí}
%
% Caja al inicio de cada sección con el presupuesto de palabras y el texto
% literal de la consigna (intro.md en el avance, formulario oficial en el
% final). Visible en borrador, invisible en FINAL.
% N = 0 significa "sin límite declarado".
% ---------------------------------------------------------------------------
\newcommand{\limite}[2]{%
  \ifdocFinal\else
    \begin{tcolorbox}[
      colback=blue!3, colframe=blue!35, boxrule=0.4pt,
      left=6pt, right=6pt, top=4pt, bottom=4pt,
      before skip=6pt, after skip=10pt]
      \footnotesize
      \ifnum#1>0\textbf{Máximo #1 palabras.}\quad\fi
      #2
    \end{tcolorbox}
  \fi
}

% ---------------------------------------------------------------------------
% \nota{...} — recordatorio interno de qué falta o de cómo se conecta con el
% resto del proyecto. Invisible en FINAL.
% ---------------------------------------------------------------------------
\newcommand{\nota}[1]{%
  \ifdocFinal\else
    {\color{orange!70!black}\footnotesize\textsf{[nota: #1]}\par\medskip}
  \fi
}

% ---------------------------------------------------------------------------
% \TODO{...} — dato que todavía no se tiene. Nunca se inventa un dato: se
% marca. En FINAL sigue VISIBLE y en rojo, a propósito: un documento no se
% entrega con campos sin llenar, y el modo final tiene que gritarlo.
% ---------------------------------------------------------------------------
\newcommand{\TODO}[1]{{\color{red}\textbf{[POR DEFINIR: #1]}}}

% common/meta-es.tex — metadatos compartidos por los documentos del curso.
% Todos en un solo sitio. Regla: aquí no se inventa nada; un dato que no se
% tiene de la fuente se deja como \TODO{...} y sale en rojo, también en FINAL.
% El nombre del documento (\docNombre) NO vive aquí: cada main-es.tex lo define,
% porque es lo único que distingue al avance del final.

% ===== Proyecto =====
\newcommand{\proyTitulo}{Planes auto-editables: Lazo abierto vs.\ replanificación}
\newcommand{\proyPregunta}{¿Bajo qué condiciones del entorno un agente que reescribe su propio plan supera a uno con plan estático y desde qué costo de reescritura se invierte esa ventaja?}

% ===== Curso =====
\newcommand{\proyCurso}{Inteligencia Artificial: Representación y Solución de Problemas}
\newcommand{\proyCursoCorto}{IA: Representación y Solución de Problemas}
\newcommand{\proyCodigo}{32310004}
\newcommand{\proyUniversidad}{Universidad del Rosario}

% ===== Gente =====
\newcommand{\proyEstudiante}{Over Haider Castrillón Valencia}
\newcommand{\proyProfesor}{Nicolás Avilán Vargas, Ph.D.}

% ===== Fecha =====
\newcommand{\proyFecha}{\today}


% --- Encabezado de página: nombre del documento (define cada main-es.tex) ---
\providecommand{\docNombre}{Documento}
\usepackage{fancyhdr}
\pagestyle{fancy}\fancyhf{}
\fancyhead[L]{\footnotesize \docNombre{} \textbar{} \proyCursoCorto}
\fancyhead[R]{\footnotesize\thepage}
\renewcommand{\headrulewidth}{0.4pt}

% --- Títulos de sección numerados, al estilo de las consignas del curso ---
\usepackage{titlesec}
\titleformat{\section}{\normalfont\large\bfseries}{\thesection.}{0.5em}{}
\titlespacing*{\section}{0pt}{1.5\baselineskip}{0.5\baselineskip}
